\subsection{Analysis of Selected SVM Kernel Configuration}
\label{sec:svm_kernel_analysis}

Hyperparameter optimization via GridSearchCV established the second-degree polynomial kernel-based SVM model as the best configuration for intersectional race and gender classification. This selected configuration applies a penalty parameter of $C=10$ without PCA or Scaler, and adopts the polynomial kernel function in \eqref{eq:8} with a scaling coefficient of $\gamma=\text{scale}$, a default constant of $\text{coef0}=0.0$, and a mapping degree of $d=2$. This empirical finding affirms that the selected configuration is the highest-performing model within the evaluated search space, rather than a claim of absolute mathematical superiority of the kernel function in general. Conceptually, the second-degree polynomial mapping enables the model to capture quadratic interactions among multi-domain latent feature dimensions without requiring an explicit infinite-dimensional projection, thereby making the resulting decision boundary more flexible in separating demographic subgroups.

Another important characteristic of the best SVM configuration in the empirical evaluation is the consistent selection of the full feature representation space without PCA dimensionality reduction (pca=None). The selection of the pca=None option indicates that retaining all 2,304 dimensions of the tri-domain latent representation provides the best classification performance within the evaluated search space. However, this interpretation is entirely grounded in empirical observational evidence, so the test results cannot be concluded as definitive proof that PCA reduction discards crucial discriminative information. The regularization characteristic of SVM through margin maximization proves effective in managing high-dimensional feature spaces without suffering performance degradation due to overfitting, thereby comprehensively preserving the latent geometric structure of the three facial representation domains.
